\chapter{Архитектура восстановления смарт-аккаунтов}\label{ch:ch5}

\section{Абстрактные аккаунты и модель зашифрованного хранения}\label{sec:ch5/sec1}

Абстракция аккаунта переносит логику аутентификации и управления средствами на
уровень смарт-контрактов, что позволяет реализовывать произвольные схемы доступа и
восстановления. В экосистеме Cosmos подобная функциональность реализуется через
смарт-контракты на базе CosmWasm и специализированные модули SDK~\cite{CosmWasm2021}. В рамках
диссертации рассматривается модель хранения зашифрованных данных с использованием
симметричного ключа, который не раскрывается в открытом виде и распределяется по
сети доверенных участников.

\subsection{Архитектурная схема и роли}\label{sec:ch5/sec1/arch}

Архитектура включает владельца аккаунта, смарт-контракт абстрактного аккаунта,
набор доверенных участников (guardians) и клиентское приложение, выполняющее
криптографические операции. Симметричный ключ, используемый для шифрования
пользовательских данных, разделяется по схеме Шамира и распределяется между
участниками~\cite{Shamir1979}. Часть долей может храниться у самого пользователя, что повышает
отказоустойчивость.

\begin{figure}[ht]
\centering
\includegraphics[width=0.8\textwidth]{scheme1.png}
\caption{Архитектура распределения симметричного ключа с использованием схемы Шамира}
\label{fig:shamir_scheme}
\end{figure}

На рис.~\ref{fig:shamir_scheme} показано, как симметричный ключ преобразуется в
набор долей, после чего каждая доля шифруется и закрепляется за конкретным
участником. Это позволяет обеспечить конфиденциальность ключа и восстановление
доступа при потере исходного секретного материала.

\subsection{Компоненты и модель данных смарт-контракта}\label{sec:ch5/sec1/storage}

Смарт-контракт хранит только зашифрованные данные и метаданные восстановления.
Ключевые элементы состояния включают:
\begin{itemize}
  \item \texttt{GUARDIANS} --- список адресов доверенных участников;
  \item \texttt{THRESHOLD} --- порог голосов для восстановления;
  \item \texttt{VOTES} --- голоса guardians по конкретному запросу восстановления;
  \item \texttt{COUNTS} --- счетчики голосов по каждому новому ключу;
  \item \texttt{KEY\_VALUE\_STORE} --- общее хранилище пользовательских данных;
  \item \texttt{DATA\_SECRET} --- зашифрованный симметричный ключ;
  \item \texttt{SHARES} --- зашифрованные доли секрета, привязанные к guardians;
  \item \texttt{RECOVER\_DATA} --- временные данные, собираемые при восстановлении.
\end{itemize}

Криптографические операции (генерация ключей, разделение секрета, шифрование долей)
выполняются на стороне клиента, что исключает раскрытие секрета в смарт-контракте и
снижает риски компрометации on-chain данных.

\section{Механизм восстановления с помощью пороговой криптографии и социального взаимодействия}\label{sec:ch5/sec2}

Для восстановления зашифрованных данных применяется схема порогового разделения
секрета Шамира. Ключ делится на части и распределяется между доверенными
участниками (guardians), при этом восстановление возможно только при достижении
порогового числа долей. Социальная модель восстановления обеспечивает баланс между
доступностью и безопасностью: компрометация менее порогового числа участников не
раскрывает секрет, а отказ части участников не блокирует процедуру восстановления.

\subsection{Процесс разделения и хранения секрета}\label{sec:ch5/sec2/share}

Процедура подготовки секрета включает следующие этапы:
\begin{enumerate}
  \item Генерация симметричного ключа на стороне клиента.
  \item Разделение ключа на $n$ долей с порогом восстановления $k$ по схеме Шамира.
  \item Шифрование каждой доли публичным ключом соответствующего guardian.
  \item Запись зашифрованных долей в состояние смарт-контракта.
\end{enumerate}

Такой подход обеспечивает конфиденциальность ключа и снижает риск его утечки даже
при частичной компрометации участников.

\subsection{Алгоритм восстановления секрета}\label{sec:ch5/sec2/recover}

Восстановление инициируется пользователем через доверенных участников и включает
голосование и сбор долей:
\begin{enumerate}
  \item Пользователь инициирует запрос на восстановление нового ключа доступа.
  \item Guardians подтверждают запрос, формируя голоса в смарт-контракте.
  \item При достижении порога \texttt{THRESHOLD} происходит обновление ключа
  аккаунта и очистка временных данных голосования.
  \item Guardians размещают зашифрованные доли секрета в \texttt{RECOVER\_DATA}.
  \item Клиентское приложение собирает достаточное число долей и восстанавливает
  симметричный ключ на стороне пользователя.
\end{enumerate}

Данный алгоритм сочетает on-chain контроль (голосование и учет порога) и off-chain
криптографические операции (восстановление секрета), что повышает безопасность и
снижает объем чувствительных данных в блокчейне.

\subsection{Клиентский интерфейс и операции}\label{sec:ch5/sec2/client}

Клиентский интерфейс предоставляет функции управления guardians, голосования и
работы с данными. На рис.~\ref{fig:account_interface} показан интерфейс управления
абстрактным аккаунтом, а на рис.~\ref{fig:secret_creation} --- процесс генерации и
распределения долей секрета.

\begin{figure}[ht]
\centering
\includegraphics[width=0.75\textwidth]{image3.png}
\caption{Интерфейс абстрактного аккаунта: управление guardians и данными}
\label{fig:account_interface}
\end{figure}

\begin{figure}[ht]
\centering
\includegraphics[width=0.75\textwidth]{image4.png}
\caption{Интерфейс создания и распределения долей секрета}
\label{fig:secret_creation}
\end{figure}

Криптографические операции выполняются на клиенте: формирование долей, их
шифрование и подпись транзакций. Такая модель сохраняет конфиденциальность ключа
и делает смарт-контракт координатором процесса, а не носителем секрета.
